\section*{AI-use statement}

This paper was drafted by a crew of AI research-assistant agents
(Free-Research-Crew, protocol v0.1) under the direction of the human author, Bryce Wang.
He approved the research question (2026-11-18, 12:24 UTC), the design (2026-11-18,
14:38 UTC) and the results (2026-11-19, 06:57 UTC) at three human gates, each in chat
after reading the Reviewer's verdict. The manuscript itself is subject to a fourth, the
submission gate. He takes responsibility for the paper's content.

\begin{itemize}
  \item \textbf{Roles and models.} Every crew role ran on Anthropic Claude through Claude
        Code. The PI role was the main Claude Code session. The Scout, Data RA, Designer,
        Econometrician, Figures RA, Reviewer and Writer ran as subagents launched from
        that session; the Replicator has not yet run. The review records name the model as
        \texttt{claude-opus-5-5}. The Econometrician's session was restarted once during
        the analysis runs (Section~\ref{sec:deviations}). Briefs to the agents were archived
        in \texttt{briefs/} from the Econometrician's relaunch brief (2026-11-18, 16:09 UTC)
        onward. Earlier briefs were not archived.
  \item \textbf{Review.} The Reviewer agent ran 4 rounds on the research question, 3 on the
        design, 2 on the code and 2 on the results. No separate literature round was held,
        and no review of the paper has been held yet. The code was reviewed, and passed,
        before any analysis run. Its independence level was \emph{fresh-context}: a new
        context in every round, with the same model family as the other roles. The human
        author chose this over a cross-family reviewer. Three final rounds passed with
        minor findings still open:
        \begin{itemize}
          \item question round~4 (three findings), carried to the design stage and resolved
                there;
          \item design round~3 (two findings), answered in the plan before the design gate;
          \item results round~2 (one finding), answered in the claims before the results
                gate.
        \end{itemize}
  \item \textbf{Findings accepted over the Reviewer's objection.} None.
  \item \textbf{Evidence.} Every result number in the text and tables is a macro generated
        from the metrics of a logged run, or is written by a logged table script from run
        metrics. The numbers typed by hand are of three kinds:
        \begin{itemize}
          \item design facts fixed in the plan before any analysis run, including the
                tolerance constants, the numbers of splits, the timing estimates, the
                half-SE thresholds, the per-cell false-failure probabilities and the
                combined 33\% rate;
          \item published or source facts, quoted as the Scout's source notes record them;
          \item process facts: dates, review-round counts, and the 92-minute wall time
                with its load figures, taken from \texttt{runs/NOTES.md} and the run
                records.
        \end{itemize}
        Every citation was checked against its source by the Scout: 17 from the full text,
        2 from the abstract only and 1 by resolving its DOI. Where only the abstract or
        metadata could be read (as for \citet{chernozhukov2004}, \citet{poterba1995} and
        \citet{abadie2003}), the paper cites the source for no more than that. The full
        workspace is public at
        \url{https://github.com/brycewang-stanford/Free-Research-Crew/tree/main/gallery/ss-dml-401k}.
        It contains the plan, the run ledger and notes, the code (with a verbatim copy of
        the authors' \texttt{ML\_Functions.R}), the data manifest, every review, the
        archived briefs and the source notes.
  \item \textbf{Replication.} Pending. No \texttt{replication.json} exists at the time of
        writing. The Replicator's clean rerun from \texttt{code/requirements.txt} is
        required before the submission gate, and its result will be reported here.
\end{itemize}
